
\documentclass{article}
\usepackage{colortbl}
\usepackage{makecell}
\usepackage{multirow}
\usepackage{supertabular}

\begin{document}

\newcounter{utterance}

\centering \large Interaction Transcript for game `codenames', experiment `opponent\_none', episode 0 with gemma{-}4{-}26b{-}a4b{-}it.
\vspace{24pt}

{ \footnotesize  \setcounter{utterance}{1}
\setlength{\tabcolsep}{0pt}
\begin{supertabular}{c@{$\;$}|p{.15\linewidth}@{}p{.15\linewidth}p{.15\linewidth}p{.15\linewidth}p{.15\linewidth}p{.15\linewidth}}
    \# & $\;$A & \multicolumn{4}{c}{Game Master} & $\;\:$B\\
    \hline

    \theutterance \stepcounter{utterance}  
    & & & \multicolumn{2}{p{0.3\linewidth}}{
        \cellcolor[rgb]{0.9,0.9,0.9}{
            \makecell[{{p{\linewidth}}}]{
                \texttt{\tiny{[GM$|$GM]}}
                {'hidden': {'team': ['data', 'kovas', 'kovotojas', 'pekinas', 'traukinys', 'varis', 'akmuo', 'greitoji pagalba', 'siena'], 'innocent': ['kamuolys', 'pirštinė', 'vanagas', 'kontrabandininkas', 'lochnesas', 'triušis', 'planšetinis kompiuteris'], 'opponent': ['karas', 'erelis', 'sniegas', 'centras', 'lenta', 'teatras', 'nykštys', 'riešutas'], 'assassin': ['mikroskopas']}, 'revealed': {'team': {'team': [], 'innocent': [], 'opponent': [], 'assassin': []}, 'opponent': {'team': [], 'innocent': [], 'opponent': [], 'assassin': []}}}
            }
        }
    }
    & & \\ \\

    \theutterance \stepcounter{utterance}  
    & & \multicolumn{4}{p{0.6\linewidth}}{
        \cellcolor[rgb]{0.9,0.9,0.9}{
            \makecell[{{p{\linewidth}}}]{
                \texttt{\tiny{[P1$\langle$GM]}}
                \texttt{Pateikite vieną vieno žodžio užuominą, susijusią su kai kuriais jūsų komandos žodžiais iš toliau pateikto sąrašo. Galite pasirinkti, su kiek ir kokiais žodžiais bus susijusi jūsų užuomina.} \\
\texttt{Tikslas {-} kad jūsų partneris, remdamasis jūsų pateiktomis užuominomis, atspėtų visus jūsų komandos žodžius, kol kita komanda neatskleidė visų savo žodžių. Jei jūsų partneris atspėja "priešininko" arba "atitraukiančio dėmesį" žodį, jo spėjimo eilė baigiasi ir priešininkų komanda gali atlikti savo ėjimą. Jei jūsų partneris atspėja "žudantį" žodį, žaidimas baigiasi ir jūsų komanda iš karto pralaimi.} \\
\texttt{Kuo daugiau žodžių, su kuriais susijusi jūsų užuomina, tuo greičiau galite laimėti žaidimą, tačiau partneriui gali būti sunkiau atspėti asociacijas, todėl išmintingai pasirinkite užuominą. Užuominos žodis turi būti semantiškai susijęs su tiksliniais žodžiais, jis negali būti vienas iš sąrašuose esančių žodžių arba sudaryti jų dalių.} \\
\\ 
\texttt{Visada pateikite vieno žodžio užuominą ir kableliais atskirtą susijusių tikslinių žodžių sąrašą toliau nurodyta forma ir atsakymus pateikite kuo trumpesnius, niekada nepateikite jokio kito teksto, nei reikalaujama šioje formoje:} \\
\\ 
\texttt{UŽUOMINA: <ŽODIS>} \\
\texttt{TIKSLAI: <TIKSLAI>} \\
\\ 
\texttt{Jūsų komandos žodžiai yra: data, kovas, kovotojas, pekinas, traukinys, varis, akmuo, greitoji pagalba, siena.} \\
\texttt{Jūsų priešininko komandos žodžiai yra: karas, erelis, sniegas, centras, lenta, teatras, nykštys, riešutas.} \\
\texttt{Atitraukimo žodžiai yra: kamuolys, pirštinė, vanagas, kontrabandininkas, lochnesas, triušis, planšetinis kompiuteris.} \\
\texttt{Žudymo žodžiai yra: mikroskopas.} \\
            }
        }
    }
    & \\ \\

    \theutterance \stepcounter{utterance}  
    & \multicolumn{4}{p{0.6\linewidth}}{
        \cellcolor[rgb]{0.9,0.9,0.9}{
            \makecell[{{p{\linewidth}}}]{
                \texttt{\tiny{[P1$\rangle$GM]}}
                \texttt{UŽUOMINA: KINETINIS} \\
\texttt{TIKSLAI: kovas, kovotojas, traukinys, akmuo} \\
            }
        }
    }
    & & \\ \\

    \theutterance \stepcounter{utterance}  
    & & & \multicolumn{2}{p{0.3\linewidth}}{
        \cellcolor[rgb]{0.9,0.9,0.9}{
            \makecell[{{p{\linewidth}}}]{
                \texttt{\tiny{[GM$|$GM]}}
                \texttt{kinetinis} \\
            }
        }
    }
    & & \\ \\

    \theutterance \stepcounter{utterance}  
    & & & \multicolumn{2}{p{0.3\linewidth}}{
        \cellcolor[rgb]{0.9,0.9,0.9}{
            \makecell[{{p{\linewidth}}}]{
                \texttt{\tiny{[GM$|$GM]}}
                ['kovas', 'kovotojas', 'traukinys', 'akmuo']
            }
        }
    }
    & & \\ \\

    \theutterance \stepcounter{utterance}  
    & & & \multicolumn{2}{p{0.3\linewidth}}{
        \cellcolor[rgb]{0.9,0.9,0.9}{
            \makecell[{{p{\linewidth}}}]{
                \texttt{\tiny{[GM$|$GM]}}
                {'word': 'kovas', 'assignment': 'team'}
            }
        }
    }
    & & \\ \\

    \theutterance \stepcounter{utterance}  
    & & & \multicolumn{2}{p{0.3\linewidth}}{
        \cellcolor[rgb]{0.9,0.9,0.9}{
            \makecell[{{p{\linewidth}}}]{
                \texttt{\tiny{[GM$|$GM]}}
                {'word': 'kovotojas', 'assignment': 'team'}
            }
        }
    }
    & & \\ \\

    \theutterance \stepcounter{utterance}  
    & & & \multicolumn{2}{p{0.3\linewidth}}{
        \cellcolor[rgb]{0.9,0.9,0.9}{
            \makecell[{{p{\linewidth}}}]{
                \texttt{\tiny{[GM$|$GM]}}
                {'word': 'traukinys', 'assignment': 'team'}
            }
        }
    }
    & & \\ \\

    \theutterance \stepcounter{utterance}  
    & & & \multicolumn{2}{p{0.3\linewidth}}{
        \cellcolor[rgb]{0.9,0.9,0.9}{
            \makecell[{{p{\linewidth}}}]{
                \texttt{\tiny{[GM$|$GM]}}
                {'word': 'akmuo', 'assignment': 'team'}
            }
        }
    }
    & & \\ \\

    \theutterance \stepcounter{utterance}  
    & & \multicolumn{4}{p{0.6\linewidth}}{
        \cellcolor[rgb]{0.9,0.9,0.9}{
            \makecell[{{p{\linewidth}}}]{
                \texttt{\tiny{[GM$\rangle$P2]}}
                \texttt{Pateikite kableliais atskirtą sąrašą iki 4 žodžių iš toliau pateikto sąrašo, kurie geriausiai siejasi arba yra labiausiai susiję su žodžiu 'kinetinis'. Visada pradėkite savo spėjimų sąrašą žodžiais SPĖJIMAS: ir į atsakymą neįtraukite jokio kito teksto.} \\
\\ 
\texttt{teatras, siena, kovas, pekinas, mikroskopas, erelis, akmuo, triušis, kovotojas, greitoji pagalba, lenta, varis, kamuolys, nykštys, riešutas, kontrabandininkas, data, centras, planšetinis kompiuteris, lochnesas, sniegas, karas, vanagas, traukinys, pirštinė} \\
            }
        }
    }
    & \\ \\

    \theutterance \stepcounter{utterance}  
    & & & \multicolumn{4}{p{0.6\linewidth}}{
        \cellcolor[rgb]{0.9,0.9,0.9}{
            \makecell[{{p{\linewidth}}}]{
                \texttt{\tiny{[GM$\langle$P2]}}
                \texttt{SPĖJIMAS: kovas, kovotojas, kamuolys, traukinys} \\
            }
        }
    }
     \\ \\

    \theutterance \stepcounter{utterance}  
    & & & \multicolumn{2}{p{0.3\linewidth}}{
        \cellcolor[rgb]{0.9,0.9,0.9}{
            \makecell[{{p{\linewidth}}}]{
                \texttt{\tiny{[GM$|$GM]}}
                ['kovas', 'kovotojas', 'kamuolys', 'traukinys']
            }
        }
    }
    & & \\ \\

    \theutterance \stepcounter{utterance}  
    & & & \multicolumn{2}{p{0.3\linewidth}}{
        \cellcolor[rgb]{0.9,0.9,0.9}{
            \makecell[{{p{\linewidth}}}]{
                \texttt{\tiny{[GM$|$GM]}}
                {'word': 'kovas', 'assignment': 'team'}
            }
        }
    }
    & & \\ \\

    \theutterance \stepcounter{utterance}  
    & & & \multicolumn{2}{p{0.3\linewidth}}{
        \cellcolor[rgb]{0.9,0.9,0.9}{
            \makecell[{{p{\linewidth}}}]{
                \texttt{\tiny{[GM$|$GM]}}
                {'word': 'kovas', 'assignment': 'team'}
            }
        }
    }
    & & \\ \\

    \theutterance \stepcounter{utterance}  
    & & & \multicolumn{2}{p{0.3\linewidth}}{
        \cellcolor[rgb]{0.9,0.9,0.9}{
            \makecell[{{p{\linewidth}}}]{
                \texttt{\tiny{[GM$|$GM]}}
                {'word': 'kovotojas', 'assignment': 'team'}
            }
        }
    }
    & & \\ \\

    \theutterance \stepcounter{utterance}  
    & & & \multicolumn{2}{p{0.3\linewidth}}{
        \cellcolor[rgb]{0.9,0.9,0.9}{
            \makecell[{{p{\linewidth}}}]{
                \texttt{\tiny{[GM$|$GM]}}
                {'word': 'kovotojas', 'assignment': 'team'}
            }
        }
    }
    & & \\ \\

    \theutterance \stepcounter{utterance}  
    & & & \multicolumn{2}{p{0.3\linewidth}}{
        \cellcolor[rgb]{0.9,0.9,0.9}{
            \makecell[{{p{\linewidth}}}]{
                \texttt{\tiny{[GM$|$GM]}}
                {'word': 'kamuolys', 'assignment': 'innocent'}
            }
        }
    }
    & & \\ \\

    \theutterance \stepcounter{utterance}  
    & & & \multicolumn{2}{p{0.3\linewidth}}{
        \cellcolor[rgb]{0.9,0.9,0.9}{
            \makecell[{{p{\linewidth}}}]{
                \texttt{\tiny{[GM$|$GM]}}
                \texttt{kamuolys} \\
            }
        }
    }
    & & \\ \\

    \theutterance \stepcounter{utterance}  
    & & & \multicolumn{2}{p{0.3\linewidth}}{
        \cellcolor[rgb]{0.9,0.9,0.9}{
            \makecell[{{p{\linewidth}}}]{
                \texttt{\tiny{[GM$|$GM]}}
                {'hidden': {'team': ['data', 'pekinas', 'traukinys', 'varis', 'akmuo', 'greitoji pagalba', 'siena'], 'innocent': ['pirštinė', 'vanagas', 'kontrabandininkas', 'lochnesas', 'triušis', 'planšetinis kompiuteris'], 'opponent': ['karas', 'erelis', 'sniegas', 'centras', 'lenta', 'teatras', 'nykštys', 'riešutas'], 'assassin': ['mikroskopas']}, 'revealed': {'team': {'team': ['kovas', 'kovotojas'], 'innocent': ['kamuolys'], 'opponent': [], 'assassin': []}, 'opponent': {'team': [], 'innocent': [], 'opponent': [], 'assassin': []}}}
            }
        }
    }
    & & \\ \\

    \theutterance \stepcounter{utterance}  
    & & \multicolumn{4}{p{0.6\linewidth}}{
        \cellcolor[rgb]{0.9,0.9,0.9}{
            \makecell[{{p{\linewidth}}}]{
                \texttt{\tiny{[P1$\langle$GM]}}
                \texttt{Žodžiai kovas, kovotojas buvo atspėti teisingai. Žodis kamuolys buvo atspėtas, tačiau tai nekaltas žodis. Jūsų komandos draugo ėjimas baigėsi.} \\
\texttt{Dabar pateikite kitą užuominą, susijusią su kai kuriais likusiais jūsų komandos žodžiais, ir susijusių tikslinių žodžių sąrašą. Nepamirškite užuominą pradėti žodžiais "UŽUOMINA: ", įrašykite naują eilutę, o tikslinių žodžių sąrašą, atskirtą kableliais, pradėkite žodžiais "TIKSLAI: ". Pastaba: palyginti su ankstesniais prašymais, kai kurie žodžiai iš sąrašų buvo pašalinti.} \\
\\ 
\texttt{Likę jūsų komandos žodžiai yra šie: data, pekinas, traukinys, varis, akmuo, greitoji pagalba, siena.} \\
\texttt{Likę priešininko žodžiai yra šie: karas, erelis, sniegas, centras, lenta, teatras, nykštys, riešutas.} \\
\texttt{Likę atitraukiantys žodžiai yra: pirštinė, vanagas, kontrabandininkas, lochnesas, triušis, planšetinis kompiuteris.} \\
\texttt{Likę nužudymo žodžiai yra: mikroskopas.} \\
            }
        }
    }
    & \\ \\

    \theutterance \stepcounter{utterance}  
    & \multicolumn{4}{p{0.6\linewidth}}{
        \cellcolor[rgb]{0.9,0.9,0.9}{
            \makecell[{{p{\linewidth}}}]{
                \texttt{\tiny{[P1$\rangle$GM]}}
                \texttt{UŽUOMINA: TRANSPORTAS} \\
\texttt{TIKSLAI: traukinys, greitoji pagalba} \\
            }
        }
    }
    & & \\ \\

    \theutterance \stepcounter{utterance}  
    & & & \multicolumn{2}{p{0.3\linewidth}}{
        \cellcolor[rgb]{0.9,0.9,0.9}{
            \makecell[{{p{\linewidth}}}]{
                \texttt{\tiny{[GM$|$GM]}}
                \texttt{transportas} \\
            }
        }
    }
    & & \\ \\

    \theutterance \stepcounter{utterance}  
    & & & \multicolumn{2}{p{0.3\linewidth}}{
        \cellcolor[rgb]{0.9,0.9,0.9}{
            \makecell[{{p{\linewidth}}}]{
                \texttt{\tiny{[GM$|$GM]}}
                ['traukinys', 'greitoji pagalba']
            }
        }
    }
    & & \\ \\

    \theutterance \stepcounter{utterance}  
    & & & \multicolumn{2}{p{0.3\linewidth}}{
        \cellcolor[rgb]{0.9,0.9,0.9}{
            \makecell[{{p{\linewidth}}}]{
                \texttt{\tiny{[GM$|$GM]}}
                {'word': 'traukinys', 'assignment': 'team'}
            }
        }
    }
    & & \\ \\

    \theutterance \stepcounter{utterance}  
    & & & \multicolumn{2}{p{0.3\linewidth}}{
        \cellcolor[rgb]{0.9,0.9,0.9}{
            \makecell[{{p{\linewidth}}}]{
                \texttt{\tiny{[GM$|$GM]}}
                {'word': 'greitoji pagalba', 'assignment': 'team'}
            }
        }
    }
    & & \\ \\

    \theutterance \stepcounter{utterance}  
    & & \multicolumn{4}{p{0.6\linewidth}}{
        \cellcolor[rgb]{0.9,0.9,0.9}{
            \makecell[{{p{\linewidth}}}]{
                \texttt{\tiny{[GM$\rangle$P2]}}
                \texttt{Žodžiai kovas, kovotojas buvo atspėti teisingai. Žodis kamuolys buvo atspėtas, tačiau tai nekaltas žodis. Jūsų ėjimas baigėsi.} \\
\texttt{Dabar pateikite dar vieną kableliais atskirtą sąrašą iš bent 1 ir ne daugiau kaip 2 žodžių iš toliau pateikto žodžių sąrašo, kurie geriausiai siejasi arba yra labiausiai susiję su žodžiu 'transportas'. Nepamirškite savo atsakymą pradėti SPĖJIMAS:. Pastaba: palyginti su ankstesniais prašymais, kai kurie žodžiai iš sąrašo buvo pašalinti.} \\
\\ 
\texttt{teatras, siena, pekinas, mikroskopas, erelis, akmuo, triušis, greitoji pagalba, lenta, varis, nykštys, riešutas, kontrabandininkas, data, centras, planšetinis kompiuteris, lochnesas, sniegas, karas, vanagas, traukinys, pirštinė} \\
            }
        }
    }
    & \\ \\

    \theutterance \stepcounter{utterance}  
    & & & \multicolumn{4}{p{0.6\linewidth}}{
        \cellcolor[rgb]{0.9,0.9,0.9}{
            \makecell[{{p{\linewidth}}}]{
                \texttt{\tiny{[GM$\langle$P2]}}
                \texttt{SPĖJIMAS: traukinys, greitoji pagalba} \\
            }
        }
    }
     \\ \\

    \theutterance \stepcounter{utterance}  
    & & & \multicolumn{2}{p{0.3\linewidth}}{
        \cellcolor[rgb]{0.9,0.9,0.9}{
            \makecell[{{p{\linewidth}}}]{
                \texttt{\tiny{[GM$|$GM]}}
                ['traukinys', 'greitoji pagalba']
            }
        }
    }
    & & \\ \\

    \theutterance \stepcounter{utterance}  
    & & & \multicolumn{2}{p{0.3\linewidth}}{
        \cellcolor[rgb]{0.9,0.9,0.9}{
            \makecell[{{p{\linewidth}}}]{
                \texttt{\tiny{[GM$|$GM]}}
                {'word': 'traukinys', 'assignment': 'team'}
            }
        }
    }
    & & \\ \\

    \theutterance \stepcounter{utterance}  
    & & & \multicolumn{2}{p{0.3\linewidth}}{
        \cellcolor[rgb]{0.9,0.9,0.9}{
            \makecell[{{p{\linewidth}}}]{
                \texttt{\tiny{[GM$|$GM]}}
                {'word': 'traukinys', 'assignment': 'team'}
            }
        }
    }
    & & \\ \\

    \theutterance \stepcounter{utterance}  
    & & & \multicolumn{2}{p{0.3\linewidth}}{
        \cellcolor[rgb]{0.9,0.9,0.9}{
            \makecell[{{p{\linewidth}}}]{
                \texttt{\tiny{[GM$|$GM]}}
                {'word': 'greitoji pagalba', 'assignment': 'team'}
            }
        }
    }
    & & \\ \\

    \theutterance \stepcounter{utterance}  
    & & & \multicolumn{2}{p{0.3\linewidth}}{
        \cellcolor[rgb]{0.9,0.9,0.9}{
            \makecell[{{p{\linewidth}}}]{
                \texttt{\tiny{[GM$|$GM]}}
                {'word': 'greitoji pagalba', 'assignment': 'team'}
            }
        }
    }
    & & \\ \\

    \theutterance \stepcounter{utterance}  
    & & & \multicolumn{2}{p{0.3\linewidth}}{
        \cellcolor[rgb]{0.9,0.9,0.9}{
            \makecell[{{p{\linewidth}}}]{
                \texttt{\tiny{[GM$|$GM]}}
                {'hidden': {'team': ['data', 'pekinas', 'varis', 'akmuo', 'siena'], 'innocent': ['pirštinė', 'vanagas', 'kontrabandininkas', 'lochnesas', 'triušis', 'planšetinis kompiuteris'], 'opponent': ['karas', 'erelis', 'sniegas', 'centras', 'lenta', 'teatras', 'nykštys', 'riešutas'], 'assassin': ['mikroskopas']}, 'revealed': {'team': {'team': ['kovas', 'kovotojas', 'traukinys', 'greitoji pagalba'], 'innocent': ['kamuolys'], 'opponent': [], 'assassin': []}, 'opponent': {'team': [], 'innocent': [], 'opponent': [], 'assassin': []}}}
            }
        }
    }
    & & \\ \\

    \theutterance \stepcounter{utterance}  
    & & \multicolumn{4}{p{0.6\linewidth}}{
        \cellcolor[rgb]{0.9,0.9,0.9}{
            \makecell[{{p{\linewidth}}}]{
                \texttt{\tiny{[P1$\langle$GM]}}
                \texttt{Žodžiai traukinys, greitoji pagalba buvo atspėti teisingai. Jūsų komandos draugo ėjimas baigėsi.} \\
\texttt{Dabar pateikite kitą užuominą, susijusią su kai kuriais likusiais jūsų komandos žodžiais, ir susijusių tikslinių žodžių sąrašą. Nepamirškite užuominą pradėti žodžiais "UŽUOMINA: ", įrašykite naują eilutę, o tikslinių žodžių sąrašą, atskirtą kableliais, pradėkite žodžiais "TIKSLAI: ". Pastaba: palyginti su ankstesniais prašymais, kai kurie žodžiai iš sąrašų buvo pašalinti.} \\
\\ 
\texttt{Likę jūsų komandos žodžiai yra šie: data, pekinas, varis, akmuo, siena.} \\
\texttt{Likę priešininko žodžiai yra šie: karas, erelis, sniegas, centras, lenta, teatras, nykštys, riešutas.} \\
\texttt{Likę atitraukiantys žodžiai yra: pirštinė, vanagas, kontrabandininkas, lochnesas, triušis, planšetinis kompiuteris.} \\
\texttt{Likę nužudymo žodžiai yra: mikroskopas.} \\
            }
        }
    }
    & \\ \\

    \theutterance \stepcounter{utterance}  
    & \multicolumn{4}{p{0.6\linewidth}}{
        \cellcolor[rgb]{0.9,0.9,0.9}{
            \makecell[{{p{\linewidth}}}]{
                \texttt{\tiny{[P1$\rangle$GM]}}
                \texttt{UŽUOMINA: METALAS} \\
\texttt{TIKSLAI: varis, akmuo} \\
            }
        }
    }
    & & \\ \\

    \theutterance \stepcounter{utterance}  
    & & & \multicolumn{2}{p{0.3\linewidth}}{
        \cellcolor[rgb]{0.9,0.9,0.9}{
            \makecell[{{p{\linewidth}}}]{
                \texttt{\tiny{[GM$|$GM]}}
                \texttt{metalas} \\
            }
        }
    }
    & & \\ \\

    \theutterance \stepcounter{utterance}  
    & & & \multicolumn{2}{p{0.3\linewidth}}{
        \cellcolor[rgb]{0.9,0.9,0.9}{
            \makecell[{{p{\linewidth}}}]{
                \texttt{\tiny{[GM$|$GM]}}
                ['varis', 'akmuo']
            }
        }
    }
    & & \\ \\

    \theutterance \stepcounter{utterance}  
    & & & \multicolumn{2}{p{0.3\linewidth}}{
        \cellcolor[rgb]{0.9,0.9,0.9}{
            \makecell[{{p{\linewidth}}}]{
                \texttt{\tiny{[GM$|$GM]}}
                {'word': 'varis', 'assignment': 'team'}
            }
        }
    }
    & & \\ \\

    \theutterance \stepcounter{utterance}  
    & & & \multicolumn{2}{p{0.3\linewidth}}{
        \cellcolor[rgb]{0.9,0.9,0.9}{
            \makecell[{{p{\linewidth}}}]{
                \texttt{\tiny{[GM$|$GM]}}
                {'word': 'akmuo', 'assignment': 'team'}
            }
        }
    }
    & & \\ \\

    \theutterance \stepcounter{utterance}  
    & & \multicolumn{4}{p{0.6\linewidth}}{
        \cellcolor[rgb]{0.9,0.9,0.9}{
            \makecell[{{p{\linewidth}}}]{
                \texttt{\tiny{[GM$\rangle$P2]}}
                \texttt{Žodžiai traukinys, greitoji pagalba buvo atspėti teisingai. Jūsų ėjimas baigėsi.} \\
\texttt{Dabar pateikite dar vieną kableliais atskirtą sąrašą iš bent 1 ir ne daugiau kaip 2 žodžių iš toliau pateikto žodžių sąrašo, kurie geriausiai siejasi arba yra labiausiai susiję su žodžiu 'metalas'. Nepamirškite savo atsakymą pradėti SPĖJIMAS:. Pastaba: palyginti su ankstesniais prašymais, kai kurie žodžiai iš sąrašo buvo pašalinti.} \\
\\ 
\texttt{teatras, siena, pekinas, mikroskopas, erelis, akmuo, triušis, lenta, varis, nykštys, riešutas, kontrabandininkas, data, centras, planšetinis kompiuteris, lochnesas, sniegas, karas, vanagas, pirštinė} \\
            }
        }
    }
    & \\ \\

    \theutterance \stepcounter{utterance}  
    & & & \multicolumn{4}{p{0.6\linewidth}}{
        \cellcolor[rgb]{0.9,0.9,0.9}{
            \makecell[{{p{\linewidth}}}]{
                \texttt{\tiny{[GM$\langle$P2]}}
                \texttt{SPĖJIMAS: varis} \\
            }
        }
    }
     \\ \\

    \theutterance \stepcounter{utterance}  
    & & & \multicolumn{2}{p{0.3\linewidth}}{
        \cellcolor[rgb]{0.9,0.9,0.9}{
            \makecell[{{p{\linewidth}}}]{
                \texttt{\tiny{[GM$|$GM]}}
                ['varis']
            }
        }
    }
    & & \\ \\

    \theutterance \stepcounter{utterance}  
    & & & \multicolumn{2}{p{0.3\linewidth}}{
        \cellcolor[rgb]{0.9,0.9,0.9}{
            \makecell[{{p{\linewidth}}}]{
                \texttt{\tiny{[GM$|$GM]}}
                {'word': 'varis', 'assignment': 'team'}
            }
        }
    }
    & & \\ \\

    \theutterance \stepcounter{utterance}  
    & & & \multicolumn{2}{p{0.3\linewidth}}{
        \cellcolor[rgb]{0.9,0.9,0.9}{
            \makecell[{{p{\linewidth}}}]{
                \texttt{\tiny{[GM$|$GM]}}
                {'word': 'varis', 'assignment': 'team'}
            }
        }
    }
    & & \\ \\

    \theutterance \stepcounter{utterance}  
    & & & \multicolumn{2}{p{0.3\linewidth}}{
        \cellcolor[rgb]{0.9,0.9,0.9}{
            \makecell[{{p{\linewidth}}}]{
                \texttt{\tiny{[GM$|$GM]}}
                {'hidden': {'team': ['data', 'pekinas', 'akmuo', 'siena'], 'innocent': ['pirštinė', 'vanagas', 'kontrabandininkas', 'lochnesas', 'triušis', 'planšetinis kompiuteris'], 'opponent': ['karas', 'erelis', 'sniegas', 'centras', 'lenta', 'teatras', 'nykštys', 'riešutas'], 'assassin': ['mikroskopas']}, 'revealed': {'team': {'team': ['kovas', 'kovotojas', 'traukinys', 'greitoji pagalba', 'varis'], 'innocent': ['kamuolys'], 'opponent': [], 'assassin': []}, 'opponent': {'team': [], 'innocent': [], 'opponent': [], 'assassin': []}}}
            }
        }
    }
    & & \\ \\

    \theutterance \stepcounter{utterance}  
    & & \multicolumn{4}{p{0.6\linewidth}}{
        \cellcolor[rgb]{0.9,0.9,0.9}{
            \makecell[{{p{\linewidth}}}]{
                \texttt{\tiny{[P1$\langle$GM]}}
                \texttt{Žodis varis buvo atspėtas teisingai. Jūsų komandos draugo ėjimas baigėsi.} \\
\texttt{Dabar pateikite kitą užuominą, susijusią su kai kuriais likusiais jūsų komandos žodžiais, ir susijusių tikslinių žodžių sąrašą. Nepamirškite užuominą pradėti žodžiais "UŽUOMINA: ", įrašykite naują eilutę, o tikslinių žodžių sąrašą, atskirtą kableliais, pradėkite žodžiais "TIKSLAI: ". Pastaba: palyginti su ankstesniais prašymais, kai kurie žodžiai iš sąrašų buvo pašalinti.} \\
\\ 
\texttt{Likę jūsų komandos žodžiai yra šie: data, pekinas, akmuo, siena.} \\
\texttt{Likę priešininko žodžiai yra šie: karas, erelis, sniegas, centras, lenta, teatras, nykštys, riešutas.} \\
\texttt{Likę atitraukiantys žodžiai yra: pirštinė, vanagas, kontrabandininkas, lochnesas, triušis, planšetinis kompiuteris.} \\
\texttt{Likę nužudymo žodžiai yra: mikroskopas.} \\
            }
        }
    }
    & \\ \\

    \theutterance \stepcounter{utterance}  
    & \multicolumn{4}{p{0.6\linewidth}}{
        \cellcolor[rgb]{0.9,0.9,0.9}{
            \makecell[{{p{\linewidth}}}]{
                \texttt{\tiny{[P1$\rangle$GM]}}
                \texttt{UŽUOMINA: KALDI} \\
\texttt{TIKSLAI: akmuo, siena} \\
            }
        }
    }
    & & \\ \\

    \theutterance \stepcounter{utterance}  
    & & & \multicolumn{2}{p{0.3\linewidth}}{
        \cellcolor[rgb]{0.9,0.9,0.9}{
            \makecell[{{p{\linewidth}}}]{
                \texttt{\tiny{[GM$|$GM]}}
                \texttt{kaldi} \\
            }
        }
    }
    & & \\ \\

    \theutterance \stepcounter{utterance}  
    & & & \multicolumn{2}{p{0.3\linewidth}}{
        \cellcolor[rgb]{0.9,0.9,0.9}{
            \makecell[{{p{\linewidth}}}]{
                \texttt{\tiny{[GM$|$GM]}}
                ['akmuo', 'siena']
            }
        }
    }
    & & \\ \\

    \theutterance \stepcounter{utterance}  
    & & & \multicolumn{2}{p{0.3\linewidth}}{
        \cellcolor[rgb]{0.9,0.9,0.9}{
            \makecell[{{p{\linewidth}}}]{
                \texttt{\tiny{[GM$|$GM]}}
                {'word': 'akmuo', 'assignment': 'team'}
            }
        }
    }
    & & \\ \\

    \theutterance \stepcounter{utterance}  
    & & & \multicolumn{2}{p{0.3\linewidth}}{
        \cellcolor[rgb]{0.9,0.9,0.9}{
            \makecell[{{p{\linewidth}}}]{
                \texttt{\tiny{[GM$|$GM]}}
                {'word': 'siena', 'assignment': 'team'}
            }
        }
    }
    & & \\ \\

    \theutterance \stepcounter{utterance}  
    & & \multicolumn{4}{p{0.6\linewidth}}{
        \cellcolor[rgb]{0.9,0.9,0.9}{
            \makecell[{{p{\linewidth}}}]{
                \texttt{\tiny{[GM$\rangle$P2]}}
                \texttt{Žodis varis buvo atspėtas teisingai. Jūsų ėjimas baigėsi.} \\
\texttt{Dabar pateikite dar vieną kableliais atskirtą sąrašą iš bent 1 ir ne daugiau kaip 2 žodžių iš toliau pateikto žodžių sąrašo, kurie geriausiai siejasi arba yra labiausiai susiję su žodžiu 'kaldi'. Nepamirškite savo atsakymą pradėti SPĖJIMAS:. Pastaba: palyginti su ankstesniais prašymais, kai kurie žodžiai iš sąrašo buvo pašalinti.} \\
\\ 
\texttt{teatras, siena, pekinas, mikroskopas, erelis, akmuo, triušis, lenta, nykštys, riešutas, kontrabandininkas, data, centras, planšetinis kompiuteris, lochnesas, sniegas, karas, vanagas, pirštinė} \\
            }
        }
    }
    & \\ \\

    \theutterance \stepcounter{utterance}  
    & & & \multicolumn{4}{p{0.6\linewidth}}{
        \cellcolor[rgb]{0.9,0.9,0.9}{
            \makecell[{{p{\linewidth}}}]{
                \texttt{\tiny{[GM$\langle$P2]}}
                \texttt{SPĖJIMAS: kontrabandininkas} \\
            }
        }
    }
     \\ \\

    \theutterance \stepcounter{utterance}  
    & & & \multicolumn{2}{p{0.3\linewidth}}{
        \cellcolor[rgb]{0.9,0.9,0.9}{
            \makecell[{{p{\linewidth}}}]{
                \texttt{\tiny{[GM$|$GM]}}
                ['kontrabandininkas']
            }
        }
    }
    & & \\ \\

    \theutterance \stepcounter{utterance}  
    & & & \multicolumn{2}{p{0.3\linewidth}}{
        \cellcolor[rgb]{0.9,0.9,0.9}{
            \makecell[{{p{\linewidth}}}]{
                \texttt{\tiny{[GM$|$GM]}}
                {'word': 'kontrabandininkas', 'assignment': 'innocent'}
            }
        }
    }
    & & \\ \\

    \theutterance \stepcounter{utterance}  
    & & & \multicolumn{2}{p{0.3\linewidth}}{
        \cellcolor[rgb]{0.9,0.9,0.9}{
            \makecell[{{p{\linewidth}}}]{
                \texttt{\tiny{[GM$|$GM]}}
                \texttt{kontrabandininkas} \\
            }
        }
    }
    & & \\ \\

    \theutterance \stepcounter{utterance}  
    & & & \multicolumn{2}{p{0.3\linewidth}}{
        \cellcolor[rgb]{0.9,0.9,0.9}{
            \makecell[{{p{\linewidth}}}]{
                \texttt{\tiny{[GM$|$GM]}}
                {'hidden': {'team': ['data', 'pekinas', 'akmuo', 'siena'], 'innocent': ['pirštinė', 'vanagas', 'lochnesas', 'triušis', 'planšetinis kompiuteris'], 'opponent': ['karas', 'erelis', 'sniegas', 'centras', 'lenta', 'teatras', 'nykštys', 'riešutas'], 'assassin': ['mikroskopas']}, 'revealed': {'team': {'team': ['kovas', 'kovotojas', 'traukinys', 'greitoji pagalba', 'varis'], 'innocent': ['kamuolys', 'kontrabandininkas'], 'opponent': [], 'assassin': []}, 'opponent': {'team': [], 'innocent': [], 'opponent': [], 'assassin': []}}}
            }
        }
    }
    & & \\ \\

    \theutterance \stepcounter{utterance}  
    & & \multicolumn{4}{p{0.6\linewidth}}{
        \cellcolor[rgb]{0.9,0.9,0.9}{
            \makecell[{{p{\linewidth}}}]{
                \texttt{\tiny{[P1$\langle$GM]}}
                \texttt{Žodis kontrabandininkas buvo atspėtas, tačiau tai nekaltas žodis. Jūsų komandos draugo ėjimas baigėsi.} \\
\texttt{Dabar pateikite kitą užuominą, susijusią su kai kuriais likusiais jūsų komandos žodžiais, ir susijusių tikslinių žodžių sąrašą. Nepamirškite užuominą pradėti žodžiais "UŽUOMINA: ", įrašykite naują eilutę, o tikslinių žodžių sąrašą, atskirtą kableliais, pradėkite žodžiais "TIKSLAI: ". Pastaba: palyginti su ankstesniais prašymais, kai kurie žodžiai iš sąrašų buvo pašalinti.} \\
\\ 
\texttt{Likę jūsų komandos žodžiai yra šie: data, pekinas, akmuo, siena.} \\
\texttt{Likę priešininko žodžiai yra šie: karas, erelis, sniegas, centras, lenta, teatras, nykštys, riešutas.} \\
\texttt{Likę atitraukiantys žodžiai yra: pirštinė, vanagas, lochnesas, triušis, planšetinis kompiuteris.} \\
\texttt{Likę nužudymo žodžiai yra: mikroskopas.} \\
            }
        }
    }
    & \\ \\

    \theutterance \stepcounter{utterance}  
    & \multicolumn{4}{p{0.6\linewidth}}{
        \cellcolor[rgb]{0.9,0.9,0.9}{
            \makecell[{{p{\linewidth}}}]{
                \texttt{\tiny{[P1$\rangle$GM]}}
                \texttt{UŽUOMINA: KALENDORIUS} \\
\texttt{TIKSLAI: data, pekinas} \\
            }
        }
    }
    & & \\ \\

    \theutterance \stepcounter{utterance}  
    & & & \multicolumn{2}{p{0.3\linewidth}}{
        \cellcolor[rgb]{0.9,0.9,0.9}{
            \makecell[{{p{\linewidth}}}]{
                \texttt{\tiny{[GM$|$GM]}}
                \texttt{kalendorius} \\
            }
        }
    }
    & & \\ \\

    \theutterance \stepcounter{utterance}  
    & & & \multicolumn{2}{p{0.3\linewidth}}{
        \cellcolor[rgb]{0.9,0.9,0.9}{
            \makecell[{{p{\linewidth}}}]{
                \texttt{\tiny{[GM$|$GM]}}
                ['data', 'pekinas']
            }
        }
    }
    & & \\ \\

    \theutterance \stepcounter{utterance}  
    & & & \multicolumn{2}{p{0.3\linewidth}}{
        \cellcolor[rgb]{0.9,0.9,0.9}{
            \makecell[{{p{\linewidth}}}]{
                \texttt{\tiny{[GM$|$GM]}}
                {'word': 'data', 'assignment': 'team'}
            }
        }
    }
    & & \\ \\

    \theutterance \stepcounter{utterance}  
    & & & \multicolumn{2}{p{0.3\linewidth}}{
        \cellcolor[rgb]{0.9,0.9,0.9}{
            \makecell[{{p{\linewidth}}}]{
                \texttt{\tiny{[GM$|$GM]}}
                {'word': 'pekinas', 'assignment': 'team'}
            }
        }
    }
    & & \\ \\

    \theutterance \stepcounter{utterance}  
    & & \multicolumn{4}{p{0.6\linewidth}}{
        \cellcolor[rgb]{0.9,0.9,0.9}{
            \makecell[{{p{\linewidth}}}]{
                \texttt{\tiny{[GM$\rangle$P2]}}
                \texttt{Žodis kontrabandininkas buvo atspėtas, tačiau tai nekaltas žodis. Jūsų ėjimas baigėsi.} \\
\texttt{Dabar pateikite dar vieną kableliais atskirtą sąrašą iš bent 1 ir ne daugiau kaip 2 žodžių iš toliau pateikto žodžių sąrašo, kurie geriausiai siejasi arba yra labiausiai susiję su žodžiu 'kalendorius'. Nepamirškite savo atsakymą pradėti SPĖJIMAS:. Pastaba: palyginti su ankstesniais prašymais, kai kurie žodžiai iš sąrašo buvo pašalinti.} \\
\\ 
\texttt{teatras, siena, pekinas, mikroskopas, erelis, akmuo, triušis, lenta, nykštys, riešutas, data, centras, planšetinis kompiuteris, lochnesas, sniegas, karas, vanagas, pirštinė} \\
            }
        }
    }
    & \\ \\

    \theutterance \stepcounter{utterance}  
    & & & \multicolumn{4}{p{0.6\linewidth}}{
        \cellcolor[rgb]{0.9,0.9,0.9}{
            \makecell[{{p{\linewidth}}}]{
                \texttt{\tiny{[GM$\langle$P2]}}
                \texttt{SPĖJIMAS: data} \\
            }
        }
    }
     \\ \\

    \theutterance \stepcounter{utterance}  
    & & & \multicolumn{2}{p{0.3\linewidth}}{
        \cellcolor[rgb]{0.9,0.9,0.9}{
            \makecell[{{p{\linewidth}}}]{
                \texttt{\tiny{[GM$|$GM]}}
                ['data']
            }
        }
    }
    & & \\ \\

    \theutterance \stepcounter{utterance}  
    & & & \multicolumn{2}{p{0.3\linewidth}}{
        \cellcolor[rgb]{0.9,0.9,0.9}{
            \makecell[{{p{\linewidth}}}]{
                \texttt{\tiny{[GM$|$GM]}}
                {'word': 'data', 'assignment': 'team'}
            }
        }
    }
    & & \\ \\

    \theutterance \stepcounter{utterance}  
    & & & \multicolumn{2}{p{0.3\linewidth}}{
        \cellcolor[rgb]{0.9,0.9,0.9}{
            \makecell[{{p{\linewidth}}}]{
                \texttt{\tiny{[GM$|$GM]}}
                {'word': 'data', 'assignment': 'team'}
            }
        }
    }
    & & \\ \\

    \theutterance \stepcounter{utterance}  
    & & & \multicolumn{2}{p{0.3\linewidth}}{
        \cellcolor[rgb]{0.9,0.9,0.9}{
            \makecell[{{p{\linewidth}}}]{
                \texttt{\tiny{[GM$|$GM]}}
                {'hidden': {'team': ['pekinas', 'akmuo', 'siena'], 'innocent': ['pirštinė', 'vanagas', 'lochnesas', 'triušis', 'planšetinis kompiuteris'], 'opponent': ['karas', 'erelis', 'sniegas', 'centras', 'lenta', 'teatras', 'nykštys', 'riešutas'], 'assassin': ['mikroskopas']}, 'revealed': {'team': {'team': ['kovas', 'kovotojas', 'traukinys', 'greitoji pagalba', 'varis', 'data'], 'innocent': ['kamuolys', 'kontrabandininkas'], 'opponent': [], 'assassin': []}, 'opponent': {'team': [], 'innocent': [], 'opponent': [], 'assassin': []}}}
            }
        }
    }
    & & \\ \\

    \theutterance \stepcounter{utterance}  
    & & \multicolumn{4}{p{0.6\linewidth}}{
        \cellcolor[rgb]{0.9,0.9,0.9}{
            \makecell[{{p{\linewidth}}}]{
                \texttt{\tiny{[P1$\langle$GM]}}
                \texttt{Žodis data buvo atspėtas teisingai. Jūsų komandos draugo ėjimas baigėsi.} \\
\texttt{Dabar pateikite kitą užuominą, susijusią su kai kuriais likusiais jūsų komandos žodžiais, ir susijusių tikslinių žodžių sąrašą. Nepamirškite užuominą pradėti žodžiais "UŽUOMINA: ", įrašykite naują eilutę, o tikslinių žodžių sąrašą, atskirtą kableliais, pradėkite žodžiais "TIKSLAI: ". Pastaba: palyginti su ankstesniais prašymais, kai kurie žodžiai iš sąrašų buvo pašalinti.} \\
\\ 
\texttt{Likę jūsų komandos žodžiai yra šie: pekinas, akmuo, siena.} \\
\texttt{Likę priešininko žodžiai yra šie: karas, erelis, sniegas, centras, lenta, teatras, nykštys, riešutas.} \\
\texttt{Likę atitraukiantys žodžiai yra: pirštinė, vanagas, lochnesas, triušis, planšetinis kompiuteris.} \\
\texttt{Likę nužudymo žodžiai yra: mikroskopas.} \\
            }
        }
    }
    & \\ \\

    \theutterance \stepcounter{utterance}  
    & \multicolumn{4}{p{0.6\linewidth}}{
        \cellcolor[rgb]{0.9,0.9,0.9}{
            \makecell[{{p{\linewidth}}}]{
                \texttt{\tiny{[P1$\rangle$GM]}}
                \texttt{UŽUOMINA: ORIENTALAS} \\
\texttt{TIKSLAI: pekinas, akmuo, siena} \\
            }
        }
    }
    & & \\ \\

    \theutterance \stepcounter{utterance}  
    & & & \multicolumn{2}{p{0.3\linewidth}}{
        \cellcolor[rgb]{0.9,0.9,0.9}{
            \makecell[{{p{\linewidth}}}]{
                \texttt{\tiny{[GM$|$GM]}}
                \texttt{orientalas} \\
            }
        }
    }
    & & \\ \\

    \theutterance \stepcounter{utterance}  
    & & & \multicolumn{2}{p{0.3\linewidth}}{
        \cellcolor[rgb]{0.9,0.9,0.9}{
            \makecell[{{p{\linewidth}}}]{
                \texttt{\tiny{[GM$|$GM]}}
                ['pekinas', 'akmuo', 'siena']
            }
        }
    }
    & & \\ \\

    \theutterance \stepcounter{utterance}  
    & & & \multicolumn{2}{p{0.3\linewidth}}{
        \cellcolor[rgb]{0.9,0.9,0.9}{
            \makecell[{{p{\linewidth}}}]{
                \texttt{\tiny{[GM$|$GM]}}
                {'word': 'pekinas', 'assignment': 'team'}
            }
        }
    }
    & & \\ \\

    \theutterance \stepcounter{utterance}  
    & & & \multicolumn{2}{p{0.3\linewidth}}{
        \cellcolor[rgb]{0.9,0.9,0.9}{
            \makecell[{{p{\linewidth}}}]{
                \texttt{\tiny{[GM$|$GM]}}
                {'word': 'akmuo', 'assignment': 'team'}
            }
        }
    }
    & & \\ \\

    \theutterance \stepcounter{utterance}  
    & & & \multicolumn{2}{p{0.3\linewidth}}{
        \cellcolor[rgb]{0.9,0.9,0.9}{
            \makecell[{{p{\linewidth}}}]{
                \texttt{\tiny{[GM$|$GM]}}
                {'word': 'siena', 'assignment': 'team'}
            }
        }
    }
    & & \\ \\

    \theutterance \stepcounter{utterance}  
    & & \multicolumn{4}{p{0.6\linewidth}}{
        \cellcolor[rgb]{0.9,0.9,0.9}{
            \makecell[{{p{\linewidth}}}]{
                \texttt{\tiny{[GM$\rangle$P2]}}
                \texttt{Žodis data buvo atspėtas teisingai. Jūsų ėjimas baigėsi.} \\
\texttt{Dabar pateikite dar vieną kableliais atskirtą sąrašą iš bent 1 ir ne daugiau kaip 3 žodžių iš toliau pateikto žodžių sąrašo, kurie geriausiai siejasi arba yra labiausiai susiję su žodžiu 'orientalas'. Nepamirškite savo atsakymą pradėti SPĖJIMAS:. Pastaba: palyginti su ankstesniais prašymais, kai kurie žodžiai iš sąrašo buvo pašalinti.} \\
\\ 
\texttt{teatras, siena, pekinas, mikroskopas, erelis, akmuo, triušis, lenta, nykštys, riešutas, centras, planšetinis kompiuteris, lochnesas, sniegas, karas, vanagas, pirštinė} \\
            }
        }
    }
    & \\ \\

    \theutterance \stepcounter{utterance}  
    & & & \multicolumn{4}{p{0.6\linewidth}}{
        \cellcolor[rgb]{0.9,0.9,0.9}{
            \makecell[{{p{\linewidth}}}]{
                \texttt{\tiny{[GM$\langle$P2]}}
                \texttt{SPĖJIMAS: pekinas} \\
            }
        }
    }
     \\ \\

    \theutterance \stepcounter{utterance}  
    & & & \multicolumn{2}{p{0.3\linewidth}}{
        \cellcolor[rgb]{0.9,0.9,0.9}{
            \makecell[{{p{\linewidth}}}]{
                \texttt{\tiny{[GM$|$GM]}}
                ['pekinas']
            }
        }
    }
    & & \\ \\

    \theutterance \stepcounter{utterance}  
    & & & \multicolumn{2}{p{0.3\linewidth}}{
        \cellcolor[rgb]{0.9,0.9,0.9}{
            \makecell[{{p{\linewidth}}}]{
                \texttt{\tiny{[GM$|$GM]}}
                {'word': 'pekinas', 'assignment': 'team'}
            }
        }
    }
    & & \\ \\

    \theutterance \stepcounter{utterance}  
    & & & \multicolumn{2}{p{0.3\linewidth}}{
        \cellcolor[rgb]{0.9,0.9,0.9}{
            \makecell[{{p{\linewidth}}}]{
                \texttt{\tiny{[GM$|$GM]}}
                {'word': 'pekinas', 'assignment': 'team'}
            }
        }
    }
    & & \\ \\

    \theutterance \stepcounter{utterance}  
    & & & \multicolumn{2}{p{0.3\linewidth}}{
        \cellcolor[rgb]{0.9,0.9,0.9}{
            \makecell[{{p{\linewidth}}}]{
                \texttt{\tiny{[GM$|$GM]}}
                {'hidden': {'team': ['akmuo', 'siena'], 'innocent': ['pirštinė', 'vanagas', 'lochnesas', 'triušis', 'planšetinis kompiuteris'], 'opponent': ['karas', 'erelis', 'sniegas', 'centras', 'lenta', 'teatras', 'nykštys', 'riešutas'], 'assassin': ['mikroskopas']}, 'revealed': {'team': {'team': ['kovas', 'kovotojas', 'traukinys', 'greitoji pagalba', 'varis', 'data', 'pekinas'], 'innocent': ['kamuolys', 'kontrabandininkas'], 'opponent': [], 'assassin': []}, 'opponent': {'team': [], 'innocent': [], 'opponent': [], 'assassin': []}}}
            }
        }
    }
    & & \\ \\

    \theutterance \stepcounter{utterance}  
    & & \multicolumn{4}{p{0.6\linewidth}}{
        \cellcolor[rgb]{0.9,0.9,0.9}{
            \makecell[{{p{\linewidth}}}]{
                \texttt{\tiny{[P1$\langle$GM]}}
                \texttt{Žodis pekinas buvo atspėtas teisingai. Jūsų komandos draugo ėjimas baigėsi.} \\
\texttt{Dabar pateikite kitą užuominą, susijusią su kai kuriais likusiais jūsų komandos žodžiais, ir susijusių tikslinių žodžių sąrašą. Nepamirškite užuominą pradėti žodžiais "UŽUOMINA: ", įrašykite naują eilutę, o tikslinių žodžių sąrašą, atskirtą kableliais, pradėkite žodžiais "TIKSLAI: ". Pastaba: palyginti su ankstesniais prašymais, kai kurie žodžiai iš sąrašų buvo pašalinti.} \\
\\ 
\texttt{Likę jūsų komandos žodžiai yra šie: akmuo, siena.} \\
\texttt{Likę priešininko žodžiai yra šie: karas, erelis, sniegas, centras, lenta, teatras, nykštys, riešutas.} \\
\texttt{Likę atitraukiantys žodžiai yra: pirštinė, vanagas, lochnesas, triušis, planšetinis kompiuteris.} \\
\texttt{Likę nužudymo žodžiai yra: mikroskopas.} \\
            }
        }
    }
    & \\ \\

    \theutterance \stepcounter{utterance}  
    & \multicolumn{4}{p{0.6\linewidth}}{
        \cellcolor[rgb]{0.9,0.9,0.9}{
            \makecell[{{p{\linewidth}}}]{
                \texttt{\tiny{[P1$\rangle$GM]}}
                \texttt{UŽUOMINA: BARIKADA} \\
\texttt{TIKSLAI: akmuo, siena} \\
            }
        }
    }
    & & \\ \\

    \theutterance \stepcounter{utterance}  
    & & & \multicolumn{2}{p{0.3\linewidth}}{
        \cellcolor[rgb]{0.9,0.9,0.9}{
            \makecell[{{p{\linewidth}}}]{
                \texttt{\tiny{[GM$|$GM]}}
                \texttt{barikada} \\
            }
        }
    }
    & & \\ \\

    \theutterance \stepcounter{utterance}  
    & & & \multicolumn{2}{p{0.3\linewidth}}{
        \cellcolor[rgb]{0.9,0.9,0.9}{
            \makecell[{{p{\linewidth}}}]{
                \texttt{\tiny{[GM$|$GM]}}
                ['akmuo', 'siena']
            }
        }
    }
    & & \\ \\

    \theutterance \stepcounter{utterance}  
    & & & \multicolumn{2}{p{0.3\linewidth}}{
        \cellcolor[rgb]{0.9,0.9,0.9}{
            \makecell[{{p{\linewidth}}}]{
                \texttt{\tiny{[GM$|$GM]}}
                {'word': 'akmuo', 'assignment': 'team'}
            }
        }
    }
    & & \\ \\

    \theutterance \stepcounter{utterance}  
    & & & \multicolumn{2}{p{0.3\linewidth}}{
        \cellcolor[rgb]{0.9,0.9,0.9}{
            \makecell[{{p{\linewidth}}}]{
                \texttt{\tiny{[GM$|$GM]}}
                {'word': 'siena', 'assignment': 'team'}
            }
        }
    }
    & & \\ \\

    \theutterance \stepcounter{utterance}  
    & & \multicolumn{4}{p{0.6\linewidth}}{
        \cellcolor[rgb]{0.9,0.9,0.9}{
            \makecell[{{p{\linewidth}}}]{
                \texttt{\tiny{[GM$\rangle$P2]}}
                \texttt{Žodis pekinas buvo atspėtas teisingai. Jūsų ėjimas baigėsi.} \\
\texttt{Dabar pateikite dar vieną kableliais atskirtą sąrašą iš bent 1 ir ne daugiau kaip 2 žodžių iš toliau pateikto žodžių sąrašo, kurie geriausiai siejasi arba yra labiausiai susiję su žodžiu 'barikada'. Nepamirškite savo atsakymą pradėti SPĖJIMAS:. Pastaba: palyginti su ankstesniais prašymais, kai kurie žodžiai iš sąrašo buvo pašalinti.} \\
\\ 
\texttt{teatras, siena, mikroskopas, erelis, akmuo, triušis, lenta, nykštys, riešutas, centras, planšetinis kompiuteris, lochnesas, sniegas, karas, vanagas, pirštinė} \\
            }
        }
    }
    & \\ \\

    \theutterance \stepcounter{utterance}  
    & & & \multicolumn{4}{p{0.6\linewidth}}{
        \cellcolor[rgb]{0.9,0.9,0.9}{
            \makecell[{{p{\linewidth}}}]{
                \texttt{\tiny{[GM$\langle$P2]}}
                \texttt{SPĖJIMAS: siena, karas} \\
            }
        }
    }
     \\ \\

    \theutterance \stepcounter{utterance}  
    & & & \multicolumn{2}{p{0.3\linewidth}}{
        \cellcolor[rgb]{0.9,0.9,0.9}{
            \makecell[{{p{\linewidth}}}]{
                \texttt{\tiny{[GM$|$GM]}}
                ['siena', 'karas']
            }
        }
    }
    & & \\ \\

    \theutterance \stepcounter{utterance}  
    & & & \multicolumn{2}{p{0.3\linewidth}}{
        \cellcolor[rgb]{0.9,0.9,0.9}{
            \makecell[{{p{\linewidth}}}]{
                \texttt{\tiny{[GM$|$GM]}}
                {'word': 'siena', 'assignment': 'team'}
            }
        }
    }
    & & \\ \\

    \theutterance \stepcounter{utterance}  
    & & & \multicolumn{2}{p{0.3\linewidth}}{
        \cellcolor[rgb]{0.9,0.9,0.9}{
            \makecell[{{p{\linewidth}}}]{
                \texttt{\tiny{[GM$|$GM]}}
                {'word': 'siena', 'assignment': 'team'}
            }
        }
    }
    & & \\ \\

    \theutterance \stepcounter{utterance}  
    & & & \multicolumn{2}{p{0.3\linewidth}}{
        \cellcolor[rgb]{0.9,0.9,0.9}{
            \makecell[{{p{\linewidth}}}]{
                \texttt{\tiny{[GM$|$GM]}}
                {'word': 'karas', 'assignment': 'opponent'}
            }
        }
    }
    & & \\ \\

    \theutterance \stepcounter{utterance}  
    & & & \multicolumn{2}{p{0.3\linewidth}}{
        \cellcolor[rgb]{0.9,0.9,0.9}{
            \makecell[{{p{\linewidth}}}]{
                \texttt{\tiny{[GM$|$GM]}}
                \texttt{karas} \\
            }
        }
    }
    & & \\ \\

    \theutterance \stepcounter{utterance}  
    & & & \multicolumn{2}{p{0.3\linewidth}}{
        \cellcolor[rgb]{0.9,0.9,0.9}{
            \makecell[{{p{\linewidth}}}]{
                \texttt{\tiny{[GM$|$GM]}}
                {'hidden': {'team': ['akmuo'], 'innocent': ['pirštinė', 'vanagas', 'lochnesas', 'triušis', 'planšetinis kompiuteris'], 'opponent': ['erelis', 'sniegas', 'centras', 'lenta', 'teatras', 'nykštys', 'riešutas'], 'assassin': ['mikroskopas']}, 'revealed': {'team': {'team': ['kovas', 'kovotojas', 'traukinys', 'greitoji pagalba', 'varis', 'data', 'pekinas', 'siena'], 'innocent': ['kamuolys', 'kontrabandininkas'], 'opponent': ['karas'], 'assassin': []}, 'opponent': {'team': [], 'innocent': [], 'opponent': [], 'assassin': []}}}
            }
        }
    }
    & & \\ \\

    \theutterance \stepcounter{utterance}  
    & & \multicolumn{4}{p{0.6\linewidth}}{
        \cellcolor[rgb]{0.9,0.9,0.9}{
            \makecell[{{p{\linewidth}}}]{
                \texttt{\tiny{[P1$\langle$GM]}}
                \texttt{Žodis siena buvo atspėtas teisingai. Žodis karas buvo atspėtas, tačiau tai priešas žodis. Jūsų komandos draugo ėjimas baigėsi.} \\
\texttt{Dabar pateikite kitą užuominą, susijusią su kai kuriais likusiais jūsų komandos žodžiais, ir susijusių tikslinių žodžių sąrašą. Nepamirškite užuominą pradėti žodžiais "UŽUOMINA: ", įrašykite naują eilutę, o tikslinių žodžių sąrašą, atskirtą kableliais, pradėkite žodžiais "TIKSLAI: ". Pastaba: palyginti su ankstesniais prašymais, kai kurie žodžiai iš sąrašų buvo pašalinti.} \\
\\ 
\texttt{Likę jūsų komandos žodžiai yra šie: akmuo.} \\
\texttt{Likę priešininko žodžiai yra šie: erelis, sniegas, centras, lenta, teatras, nykštys, riešutas.} \\
\texttt{Likę atitraukiantys žodžiai yra: pirštinė, vanagas, lochnesas, triušis, planšetinis kompiuteris.} \\
\texttt{Likę nužudymo žodžiai yra: mikroskopas.} \\
            }
        }
    }
    & \\ \\

    \theutterance \stepcounter{utterance}  
    & \multicolumn{4}{p{0.6\linewidth}}{
        \cellcolor[rgb]{0.9,0.9,0.9}{
            \makecell[{{p{\linewidth}}}]{
                \texttt{\tiny{[P1$\rangle$GM]}}
                \texttt{UŽUOMINA: KLASTA} \\
\texttt{TIKSLAI: akmuo} \\
            }
        }
    }
    & & \\ \\

    \theutterance \stepcounter{utterance}  
    & & & \multicolumn{2}{p{0.3\linewidth}}{
        \cellcolor[rgb]{0.9,0.9,0.9}{
            \makecell[{{p{\linewidth}}}]{
                \texttt{\tiny{[GM$|$GM]}}
                \texttt{klasta} \\
            }
        }
    }
    & & \\ \\

    \theutterance \stepcounter{utterance}  
    & & & \multicolumn{2}{p{0.3\linewidth}}{
        \cellcolor[rgb]{0.9,0.9,0.9}{
            \makecell[{{p{\linewidth}}}]{
                \texttt{\tiny{[GM$|$GM]}}
                ['akmuo']
            }
        }
    }
    & & \\ \\

    \theutterance \stepcounter{utterance}  
    & & & \multicolumn{2}{p{0.3\linewidth}}{
        \cellcolor[rgb]{0.9,0.9,0.9}{
            \makecell[{{p{\linewidth}}}]{
                \texttt{\tiny{[GM$|$GM]}}
                {'word': 'akmuo', 'assignment': 'team'}
            }
        }
    }
    & & \\ \\

    \theutterance \stepcounter{utterance}  
    & & \multicolumn{4}{p{0.6\linewidth}}{
        \cellcolor[rgb]{0.9,0.9,0.9}{
            \makecell[{{p{\linewidth}}}]{
                \texttt{\tiny{[GM$\rangle$P2]}}
                \texttt{Žodis siena buvo atspėtas teisingai. Žodis karas buvo atspėtas, tačiau tai priešas žodis. Jūsų ėjimas baigėsi.} \\
\texttt{Dabar pateikite dar vieną kableliais atskirtą sąrašą iš bent 1 ir ne daugiau kaip 1 žodžių iš toliau pateikto žodžių sąrašo, kurie geriausiai siejasi arba yra labiausiai susiję su žodžiu 'klasta'. Nepamirškite savo atsakymą pradėti SPĖJIMAS:. Pastaba: palyginti su ankstesniais prašymais, kai kurie žodžiai iš sąrašo buvo pašalinti.} \\
\\ 
\texttt{teatras, mikroskopas, erelis, akmuo, triušis, lenta, nykštys, riešutas, centras, planšetinis kompiuteris, lochnesas, sniegas, vanagas, pirštinė} \\
            }
        }
    }
    & \\ \\

    \theutterance \stepcounter{utterance}  
    & & & \multicolumn{4}{p{0.6\linewidth}}{
        \cellcolor[rgb]{0.9,0.9,0.9}{
            \makecell[{{p{\linewidth}}}]{
                \texttt{\tiny{[GM$\langle$P2]}}
                \texttt{SPĖJIMAS: triušis} \\
            }
        }
    }
     \\ \\

    \theutterance \stepcounter{utterance}  
    & & & \multicolumn{2}{p{0.3\linewidth}}{
        \cellcolor[rgb]{0.9,0.9,0.9}{
            \makecell[{{p{\linewidth}}}]{
                \texttt{\tiny{[GM$|$GM]}}
                ['triušis']
            }
        }
    }
    & & \\ \\

    \theutterance \stepcounter{utterance}  
    & & & \multicolumn{2}{p{0.3\linewidth}}{
        \cellcolor[rgb]{0.9,0.9,0.9}{
            \makecell[{{p{\linewidth}}}]{
                \texttt{\tiny{[GM$|$GM]}}
                {'word': 'triušis', 'assignment': 'innocent'}
            }
        }
    }
    & & \\ \\

    \theutterance \stepcounter{utterance}  
    & & & \multicolumn{2}{p{0.3\linewidth}}{
        \cellcolor[rgb]{0.9,0.9,0.9}{
            \makecell[{{p{\linewidth}}}]{
                \texttt{\tiny{[GM$|$GM]}}
                \texttt{triušis} \\
            }
        }
    }
    & & \\ \\

    \theutterance \stepcounter{utterance}  
    & & & \multicolumn{2}{p{0.3\linewidth}}{
        \cellcolor[rgb]{0.9,0.9,0.9}{
            \makecell[{{p{\linewidth}}}]{
                \texttt{\tiny{[GM$|$GM]}}
                {'hidden': {'team': ['akmuo'], 'innocent': ['pirštinė', 'vanagas', 'lochnesas', 'planšetinis kompiuteris'], 'opponent': ['erelis', 'sniegas', 'centras', 'lenta', 'teatras', 'nykštys', 'riešutas'], 'assassin': ['mikroskopas']}, 'revealed': {'team': {'team': ['kovas', 'kovotojas', 'traukinys', 'greitoji pagalba', 'varis', 'data', 'pekinas', 'siena'], 'innocent': ['kamuolys', 'kontrabandininkas', 'triušis'], 'opponent': ['karas'], 'assassin': []}, 'opponent': {'team': [], 'innocent': [], 'opponent': [], 'assassin': []}}}
            }
        }
    }
    & & \\ \\

    \theutterance \stepcounter{utterance}  
    & & \multicolumn{4}{p{0.6\linewidth}}{
        \cellcolor[rgb]{0.9,0.9,0.9}{
            \makecell[{{p{\linewidth}}}]{
                \texttt{\tiny{[P1$\langle$GM]}}
                \texttt{Žodis triušis buvo atspėtas, tačiau tai nekaltas žodis. Jūsų komandos draugo ėjimas baigėsi.} \\
\texttt{Dabar pateikite kitą užuominą, susijusią su kai kuriais likusiais jūsų komandos žodžiais, ir susijusių tikslinių žodžių sąrašą. Nepamirškite užuominą pradėti žodžiais "UŽUOMINA: ", įrašykite naują eilutę, o tikslinių žodžių sąrašą, atskirtą kableliais, pradėkite žodžiais "TIKSLAI: ". Pastaba: palyginti su ankstesniais prašymais, kai kurie žodžiai iš sąrašų buvo pašalinti.} \\
\\ 
\texttt{Likę jūsų komandos žodžiai yra šie: akmuo.} \\
\texttt{Likę priešininko žodžiai yra šie: erelis, sniegas, centras, lenta, teatras, nykštys, riešutas.} \\
\texttt{Likę atitraukiantys žodžiai yra: pirštinė, vanagas, lochnesas, planšetinis kompiuteris.} \\
\texttt{Likę nužudymo žodžiai yra: mikroskopas.} \\
            }
        }
    }
    & \\ \\

    \theutterance \stepcounter{utterance}  
    & \multicolumn{4}{p{0.6\linewidth}}{
        \cellcolor[rgb]{0.9,0.9,0.9}{
            \makecell[{{p{\linewidth}}}]{
                \texttt{\tiny{[P1$\rangle$GM]}}
                \texttt{UŽUOMINA: Kietas} \\
\texttt{TIKSLAI: akmuo} \\
            }
        }
    }
    & & \\ \\

    \theutterance \stepcounter{utterance}  
    & & & \multicolumn{2}{p{0.3\linewidth}}{
        \cellcolor[rgb]{0.9,0.9,0.9}{
            \makecell[{{p{\linewidth}}}]{
                \texttt{\tiny{[GM$|$GM]}}
                \texttt{kietas} \\
            }
        }
    }
    & & \\ \\

    \theutterance \stepcounter{utterance}  
    & & & \multicolumn{2}{p{0.3\linewidth}}{
        \cellcolor[rgb]{0.9,0.9,0.9}{
            \makecell[{{p{\linewidth}}}]{
                \texttt{\tiny{[GM$|$GM]}}
                ['akmuo']
            }
        }
    }
    & & \\ \\

    \theutterance \stepcounter{utterance}  
    & & & \multicolumn{2}{p{0.3\linewidth}}{
        \cellcolor[rgb]{0.9,0.9,0.9}{
            \makecell[{{p{\linewidth}}}]{
                \texttt{\tiny{[GM$|$GM]}}
                {'word': 'akmuo', 'assignment': 'team'}
            }
        }
    }
    & & \\ \\

    \theutterance \stepcounter{utterance}  
    & & \multicolumn{4}{p{0.6\linewidth}}{
        \cellcolor[rgb]{0.9,0.9,0.9}{
            \makecell[{{p{\linewidth}}}]{
                \texttt{\tiny{[GM$\rangle$P2]}}
                \texttt{Žodis triušis buvo atspėtas, tačiau tai nekaltas žodis. Jūsų ėjimas baigėsi.} \\
\texttt{Dabar pateikite dar vieną kableliais atskirtą sąrašą iš bent 1 ir ne daugiau kaip 1 žodžių iš toliau pateikto žodžių sąrašo, kurie geriausiai siejasi arba yra labiausiai susiję su žodžiu 'kietas'. Nepamirškite savo atsakymą pradėti SPĖJIMAS:. Pastaba: palyginti su ankstesniais prašymais, kai kurie žodžiai iš sąrašo buvo pašalinti.} \\
\\ 
\texttt{teatras, mikroskopas, erelis, akmuo, lenta, nykštys, riešutas, centras, planšetinis kompiuteris, lochnesas, sniegas, vanagas, pirštinė} \\
            }
        }
    }
    & \\ \\

    \theutterance \stepcounter{utterance}  
    & & & \multicolumn{4}{p{0.6\linewidth}}{
        \cellcolor[rgb]{0.9,0.9,0.9}{
            \makecell[{{p{\linewidth}}}]{
                \texttt{\tiny{[GM$\langle$P2]}}
                \texttt{SPĖJIMAS: akmuo} \\
            }
        }
    }
     \\ \\

    \theutterance \stepcounter{utterance}  
    & & & \multicolumn{2}{p{0.3\linewidth}}{
        \cellcolor[rgb]{0.9,0.9,0.9}{
            \makecell[{{p{\linewidth}}}]{
                \texttt{\tiny{[GM$|$GM]}}
                ['akmuo']
            }
        }
    }
    & & \\ \\

    \theutterance \stepcounter{utterance}  
    & & & \multicolumn{2}{p{0.3\linewidth}}{
        \cellcolor[rgb]{0.9,0.9,0.9}{
            \makecell[{{p{\linewidth}}}]{
                \texttt{\tiny{[GM$|$GM]}}
                {'word': 'akmuo', 'assignment': 'team'}
            }
        }
    }
    & & \\ \\

    \theutterance \stepcounter{utterance}  
    & & & \multicolumn{2}{p{0.3\linewidth}}{
        \cellcolor[rgb]{0.9,0.9,0.9}{
            \makecell[{{p{\linewidth}}}]{
                \texttt{\tiny{[GM$|$GM]}}
                {'word': 'akmuo', 'assignment': 'team'}
            }
        }
    }
    & & \\ \\

    \theutterance \stepcounter{utterance}  
    & & & \multicolumn{2}{p{0.3\linewidth}}{
        \cellcolor[rgb]{0.9,0.9,0.9}{
            \makecell[{{p{\linewidth}}}]{
                \texttt{\tiny{[GM$|$GM]}}
                \texttt{team has won} \\
            }
        }
    }
    & & \\ \\

\end{supertabular}
}

\end{document}
